%% ma: L12-B14-0022
%% lop: 12
%% chuong: 5
%% bai: 14
%% dang: 12.14.17
%% loai: TN4
%% muc: TH
%% dap_an: "C"
%% nguon: "CĐ2.5.pdf giáo viên gửi 27/10/2026 (M07, Lê Minh Khôi – Nguyễn Anh Khôi, Chương II Chủ đề 5 Tích có hướng của hai vectơ), Câu 6"
%% kien_thuc: "Ba vectơ đồng phẳng, tích hỗn tạp (chương trình cũ, không có trong SGK KNTT 12)"
%% trang_thai: cach-ly
%% ly_do: "Vượt chương trình: điều kiện đồng phẳng của ba vectơ (tích hỗn tạp) không có trong SGK KNTT 12. Gợi ý: bỏ, hoặc đổi thành câu bốn điểm đồng phẳng giải bằng phương trình mặt phẳng."
\begin{ex}
Trong không gian $Oxyz$, điều kiện để ba vectơ $\vec a$, $\vec b$, $\vec c$ đồng phẳng là
\choice{$\left[\left[\vec a,\vec b\right],\vec c\right]=\vec 0$}{$\left[\vec a,\vec b\right]\cdot\vec c=\vec 0$}{\True $\left[\vec a,\vec b\right]\cdot\vec c=0$}{$\left[\left[\vec a,\vec b\right],\vec c\right]=0$}
\loigiai{Ba vectơ đồng phẳng $\Leftrightarrow\left[\vec a,\vec b\right]\cdot\vec c=0$ (tích vô hướng là một số).}
\end{ex}
